%\clearpage
%\chapter*{Abstract}
%\addcontentsline{toc}{chapter}{\protect Abstract}

In this thesis, we investigate how structural disorder affects ultrafast electronic dynamics probed in attosecond pump–probe experiments. Using real-time time-dependent density functional theory (RT-TDDFT) simulations performed with the SALMON code, we analyze disorder-induced modifications of carrier photoinjection, intraband current relaxation, and coherent current oscillations in semiconducting silicon.

We demonstrate that disorder reshapes the density of states, reduces the band gap, and introduces elastic scattering that governs momentum relaxation. In an orthogonal pump–probe configuration, the probe-induced intraband current is described within a generalized Drude framework, enabling extraction of the transport relaxation time. The analysis reveals that the effective mass is not constant but depends on the crystal momentum, reflecting the underlying band curvature and its disorder-induced modification.

In the field-free regime, coherent current oscillations are analyzed using an autocorrelation-based approach, allowing a decomposition of dephasing into homogeneous (scattering-induced) and inhomogeneous (dispersion-induced) contributions. A central result is the quantitative link between the homogeneous dephasing time and the transport relaxation time, indicating that the same elastic scattering mechanism controls both intraband transport and interband coherence decay.
